\section{Methodology}

This project follows an iterative, modular development methodology combining full-stack web development with experimental 3D reconstruction workflows. The system was developed in two parallel components: (1) a digital cultural archive platform and (2) a 3D scene reconstruction module. The workflow integrates data collection, preprocessing, model generation, and deployment into a unified pipeline.

\subsection{System Development Approach}

The development process followed an incremental model, where core archival functionality was implemented first, followed by integration of advanced visualization features. The major phases included:

\begin{enumerate}
    \item Requirements gathering and domain understanding
    \item System architecture and database design
    \item Implementation of media archival platform
    \item Integration of transcription and search features
    \item Experimentation with 3D reconstruction techniques
    \item Transition to 3D Gaussian Splatting for real-time rendering
    \item Deployment and performance optimization
\end{enumerate}

\subsection{Data Collection}

Cultural data was collected in the form of:
\begin{itemize}
    \item Photographs of community environments
    \item Audio interviews and oral histories
    \item Video recordings of cultural practices
    \item Contextual metadata and consent documentation
\end{itemize}

For 3D reconstruction, multi-view images were captured with overlapping viewpoints to ensure sufficient coverage of the scene. Typically, between 20 and 200 images were used for each reconstruction.

\subsection{Preprocessing Pipeline}

Before reconstruction, the collected images underwent preprocessing:

\begin{itemize}
    \item Image resizing and format normalization
    \item Removal of blurry or redundant frames
    \item EXIF extraction for camera information (when available)
    \item Structure-from-Motion (SfM) using COLMAP for camera pose estimation
\end{itemize}

The output of this stage included camera intrinsics, extrinsics, and a sparse point cloud used for initialization.

\subsection{3D Reconstruction Method}

Initially, Neural Radiance Fields (NeRF) were explored for scene reconstruction. While NeRF produced high-quality outputs, its slow training time and rendering latency made it unsuitable for real-time web-based deployment. Therefore, the methodology transitioned to 3D Gaussian Splatting (3DGS) as the primary reconstruction approach.

In 3DGS, the scene is represented as a set of Gaussian primitives. The reconstruction process includes:

\begin{enumerate}
    \item Initializing Gaussians from sparse point cloud
    \item Optimizing Gaussian parameters (position, color, covariance, opacity)
    \item Minimizing photometric reconstruction error
    \item Generating real-time renderable scene representation
\end{enumerate}

This method significantly reduced training time while enabling interactive visualization.

\subsection{Web Integration}

The reconstructed scenes were integrated into the web platform using a browser-based 3D viewer. The integration involved:

\begin{itemize}
    \item Exporting optimized Gaussian representation
    \item Converting assets to web-compatible format
    \item Loading scene via Three.js-based renderer
    \item Linking scenes to community pages in the archive
\end{itemize}

Users can navigate reconstructed environments interactively, allowing immersive exploration of cultural spaces.

\subsection{Search and Metadata Indexing}

All uploaded media and transcripts were indexed using a search engine. The indexing pipeline included:

\begin{itemize}
    \item Metadata extraction
    \item Transcript generation using speech-to-text
    \item Full-text indexing
    \item Tag-based filtering
\end{itemize}

This enabled efficient retrieval of cultural artifacts and contextual information.

\subsection{Evaluation Criteria}

The system was evaluated based on:

\begin{itemize}
    \item Reconstruction quality
    \item Rendering performance
    \item Scene load time
    \item Search accuracy
    \item Usability of archival interface
\end{itemize}

3D Gaussian Splatting demonstrated improved performance compared to NeRF, particularly in rendering speed and interactivity, making it suitable for deployment in a web-based cultural archive.